\centerline{\bf \Huge Геометрия II}
\centerline{\bf \Huge Удвоение медианы}

\begin{flushright}
        \scriptsize{Одна медиана хорошо, а две лучше}
\end{flushright}

\problem{1}
Докажите равенство треугольников по двум сторонам и медиане, исходящим из одной вершины.

\problem{2}
Докажите равенство треугольников по медиане и углам, на которые медиана разбивает угол треугольника.

\problem{3}
На стороне $A C$ треугольника $A B C$ нашлись такие точки $K$ и $L$, что $L-$ середина $A K$ и $B K$ - биссектриса угла $L B C$. Оказалось, что $B C=2 B L$. Докажите, что $K C=A B$.

\problem{4}
На стороне $B C$ треугольника $A B C$ выбрана точка $F$. Оказалось, что отрезок $A F$ пересекает медиану $B D$ в точке $E$ так, что $A E=B C$. Докажите, что $B F=F E$.

\problem{5}
На двух сторонах треугольника вне его построены квадраты. Докажите, что отрезок, соединяющий концы сторон квадратов, выходящих из одной вершины треугольника, в два раза больше медианы треугольника, выходящей из той же вершины.

\problem{6}
Точка $M$ - середина стороны $B C$ треугольника $A B C$. Из вершины $C$ опущен перпендикуляр $C L$ на прямую $A M$ ( $L$ лежит между $A$ и $M$ ). На отрезке $A M$ отмечена точка $K$ так, что $A K=2 L M$. Докажите, что $\angle B K M=\angle C A M$.

 Пусть $ABC$ --- неравнобедренный треугольник с ортоцентром $H$ и описанной окружностью $\Omega$. Прямая, проходящая через $H$, пересекает отрезки $AB$ и $AC$ в точках $E$ и $F$ соответственно. Пусть $K$ --- центр окружности $(AEF)$, и пусть прямая $AK$ вторично пересекает $\Omega$ в точке $D$. Докажите, что прямая $HK$ и прямая, проходящая через $D$ перпендикулярно $BC$, пересекаются на $\Omega$.

